\documentclass[11pt, a4paper]{article}

% Paquetes esenciales
\usepackage[utf8]{inputenc}
\usepackage[spanish, es-tabla]{babel}
\usepackage[margin=2.5cm]{geometry}
\usepackage{amsmath, amssymb}
\usepackage{graphicx}
\usepackage{xcolor}
\usepackage{tcolorbox}
\usepackage{booktabs}
\usepackage{enumitem}
\usepackage{fancyhdr}

% Paquetes para esquemas vectoriales
\usepackage{tikz}
% arrows.meta: necesario para las puntas de flecha tipo 'Stealth'
% positioning: necesario para usar 'right=of'
% babel: soluciona automáticamente los conflictos con babel-spanish y caracteres activos (<, >)
\usetikzlibrary{positioning, arrows.meta, babel}

% Configuración de hipervínculos (hyperref SIEMPRE debe ser el último paquete)
\usepackage{hyperref}
\hypersetup{
    colorlinks=true,
    linkcolor=blue!70!black,
    urlcolor=blue!70!black,
    pdftitle={Guía PFI - Circuitos Electrónicos III}
}

% Configuración de encabezados y pies de página
\pagestyle{fancy}
% Aumentado a 15pt para evitar la advertencia de compilación de fancyhdr con fuentes 11pt
\setlength{\headheight}{15pt}
\fancyhf{}
\fancyhead[L]{\textbf{Circuitos Electrónicos III (IMT-221)}}
\fancyhead[R]{Proyecto Final Integrador (PFI)}
\fancyfoot[C]{\thepage}
\fancyfoot[R]{\footnotesize Semestre 2-2026}
\renewcommand{\headrulewidth}{0.4pt}
\renewcommand{\footrulewidth}{0.4pt}

% Entornos personalizados
\newtcolorbox{importante}[1][]{
    colback=red!5!white,
    colframe=red!75!black,
    fonttitle=\bfseries,
    title=#1
}

\newtcolorbox{info}[1][]{
    colback=blue!5!white,
    colframe=blue!75!black,
    fonttitle=\bfseries,
    title=#1
}

\begin{document}

% ==========================================
% PORTADA / ENCABEZADO
% ==========================================
\begin{center}
    {\Large \textbf{UNIVERSIDAD CATÓLICA BOLIVIANA ``SAN PABLO'' – SEDE TARIJA}}\\[0.2cm]
    {\large Departamento de Ciencias de la Tecnología e Innovación}\\[0.5cm]
    
    \rule{\textwidth}{0.5pt}\\[0.4cm]
    {\huge \textbf{DOCUMENTO MAESTRO: PITCH Y GUÍA DE INICIO}}\\[0.2cm]
    {\Large \textbf{PROYECTO FINAL INTEGRADOR (PFI)}}\\[0.2cm]
    \rule{\textwidth}{0.5pt}\\[0.5cm]
\end{center}

\begin{flushleft}
    \textbf{Carrera:} Ingeniería Mecatrónica / Ingeniería Biomédica \\
    \textbf{Asignatura:} Circuitos Electrónicos III (IMT-221) – Semestre 2-2026 \\
    \textbf{Docente:} M.Sc. Ing. Bernardo Quiroga Turdera
\end{flushleft}

\vspace{0.5cm}

% ==========================================
% 1. EL RETO DE INGENIERÍA
% ==========================================
\section{El Reto de Ingeniería: ¿Por qué este proyecto?}

En el mundo profesional actual, los microcontroladores y computadoras no pueden conectarse directamente con el mundo físico. Los sensores entregan señales extremadamente débiles e inmersas en ruido, mientras que los actuadores industriales o médicos demandan altas corrientes y voltajes con dinámicas rápidas. La electrónica analógica no es una disciplina del pasado; es la columna vertebral física que permite a un sistema embebido medir con precisión y actuar sin destruirse.

En esta asignatura no construiremos circuitos de juguete ni memorizaremos fórmulas aisladas. Desde el primer día de clases, asumirás el rol de un ingeniero de desarrollo de hardware analógico, diseñando e implementando un sistema de potencia, sensado y control de lazo cerrado para un entorno real.

% ==========================================
% 2. DEFINICIÓN Y CONTEXTUALIZACIÓN DUAL
% ==========================================
\section{Definición y Contextualización Dual del Proyecto}

El Proyecto Final Integrador (PFI) consiste en el \textbf{``Diseño, Simulación e Implementación de una Etapa Analógica Bidireccional de Potencia, Medición de Corriente y Protección''}. Para responder a tu perfil profesional específico, el proyecto adopta una aplicación diferenciada según tu carrera:

\begin{info}[Para Estudiantes de Ingeniería Mecatrónica]
    \textbf{Título del Proyecto:} Driver Inteligente de Lazo Cerrado para el Control de Velocidad y Torque de un Actuador DC.\\
    \textbf{Aplicación:} Robótica móvil, servocontroladores industriales y posicionamiento de precisión.
\end{info}

\begin{info}[Para Estudiantes de Ingeniería Biomédica]
    \textbf{Título del Proyecto:} Etapa de Control de Presión y Sensado Analógico de Corriente para Detección de Oclusiones en Bombas de Infusión Clínica.\\
    \textbf{Aplicación:} Equipos de soporte de vida y bombas volumétricas. Si la vía de fluidos de un paciente se obstruye, la corriente del motor se eleva de inmediato; tu circuito analógico detectará este cambio de milivoltios de forma instantánea y protegerá el sistema.
\end{info}

% ==========================================
% 3. ARQUITECTURA DEL SISTEMA Y LOS 4 HITOS
% ==========================================
\section{Arquitectura del Sistema y los 4 Hitos de Madurez Técnica}

El proyecto no se entrega al final del semestre; se construye de forma incremental a lo largo de 4 Hitos evaluativos, donde cada unidad temática aporta un bloque físico funcional:

\begin{center}
\resizebox{\textwidth}{!}{% <-- Ajusta todo el gráfico al ancho exacto de los márgenes
\begin{tikzpicture}[
    node distance=1.5cm, % Reducido ligeramente para mejor proporción
    hito/.style={rectangle, draw=blue!80!black, fill=blue!10, text width=3.5cm, align=center, rounded corners, minimum height=1.5cm, font=\small},
    flecha/.style={thick, ->, >=Stealth}
]
    \node[hito] (h1) {\textbf{Hito 1: Sem 4}\\Protección y Clamping};
    \node[hito, right=of h1] (h2) {\textbf{Hito 2: Sem 8}\\Sensado por Shunt};
    \node[hito, right=of h2] (h3) {\textbf{Hito 3: Sem 13}\\Puente H de Potencia};
    \node[hito, right=of h3] (h4) {\textbf{Hito 4: Sem 16}\\Lazo Cerrado y Estabilidad};

    \draw[flecha] (h1) -- (h2);
    \draw[flecha] (h2) -- (h3);
    \draw[flecha] (h3) -- (h4);
\end{tikzpicture}%
}
\end{center}

\begin{description}[leftmargin=3.2cm]
    \item[Hito 1 (Semana 4):] \textbf{Módulo de Protección y Clamping Analógico.}\\
    Diseño e implementación de diodos Flyback para absorber transitorios inductivos y una red limitadora Zener ($0\text{ V}$ a $+5.1\text{ V}$) que protege el conversor analógico-digital (ADC).
    
    \item[Hito 2 (Semana 8):] \textbf{Módulo de Sensado de Corriente de Precisión.}\\
    Diseño de un amplificador diferencial discreto con transistores BJT para medir la caída de milivoltios en una resistencia shunt ($0.1\,\Omega$), eliminando el ruido eléctrico ($CMRR \ge 60\text{ dB}$).
    
    \item[Hito 3 (Semana 13):] \textbf{Módulo de Actuación y Potencia.}\\
    Diseño de un Puente H discreto con MOSFETs complementarios (canal N y P) con sus respectivos excitadores de compuerta (Gate Drivers) discretos y disipadores térmicos dimensionados matemáticamente.
    
    \item[Hito 4 (Semana 16):] \textbf{Lazo Cerrado y Estabilidad Dinámica.}\\
    Diseño de un compensador analógico Lead-Lag para cerrar el lazo de corriente, garantizando que el sistema sea estable ($PM \ge 60^\circ$) ante variaciones de carga.
\end{description}

% ==========================================
% 4. METODOLOGÍA DE TRABAJO
% ==========================================
\section{Metodología de Trabajo: Los 4 Pilares Prácticos y Reglas del Juego}

Cada entregable del proyecto se rige bajo la metodología internacional \textbf{CDIO (Concebir - Diseñar - Implementar - Operar)}. Para aprobar cada hito, tu equipo debe demostrar cuatro pilares:

\begin{enumerate}
    \item \textbf{Fundamentación Teórica en Python:} No se acepta el diseño por ensayo y error. Cada resistor, capacitor y transistor debe estar calculado mediante modelos matemáticos escritos en scripts interactivos de Python (Jupyter Notebooks).
    \item \textbf{Simulación en LTSpice:} Validación de las formas de onda teóricas antes de tocar el hardware.
    \item \textbf{Implementación Física de Hardware:} Montaje ordenado, limpio y seguro del circuito en tarjeta de prototipos.
    \item \textbf{Validación con Instrumentación Real:} Adquisición de señales en vivo mediante osciloscopio digital. Los archivos de datos \texttt{.csv} reales capturados en el laboratorio deberán importarse en Python para graficar la superposición entre la Teoría, la Simulación y la Práctica.
\end{enumerate}

\begin{importante}[Aviso de Control de Inteligencia Artificial (Política de Auditoría Dinámica)]
    Puedes usar la IA como copiloto para generar código inicial. Sin embargo, la evaluación es \textbf{individual y en vivo en tu mesa de trabajo}. El docente te pedirá modificar variables en tu código de Python en caliente (por ejemplo, cambiar la frecuencia o una capacitancia) y deberás predecir verbalmente y demostrar en el osciloscopio el impacto físico en el circuito. \textbf{Si no puedes defender tu código, el hito no será aprobado.}
\end{importante}

% ==========================================
% 5. ENTREGABLE FINAL
% ==========================================
\section{Entregable Final y Evaluación (Semana 18)}

El semestre culmina en la Semana 18 con dos entregables integradores:

\begin{itemize}
    \item \textbf{Demostración del Hardware en Vivo:} Presentación pública de tu placa funcionando de forma estable ante un panel evaluador de ingenieros.
    \item \textbf{Artículo Técnico Corto (Formato IEEE):} Un reporte científico de 4 páginas que documenta el modelado, las simulaciones, el cálculo térmico y la validación de estabilidad de tu sistema analógico.
\end{itemize}

\end{document}